\documentclass[11pt,a4paper]{article}
\usepackage[top=2.3cm,bottom=2.3cm,left=2.2cm,right=2.2cm]{geometry}

\usepackage{amsmath,amssymb}
\usepackage{fontspec}
\setmainfont{Liberation Serif}
\setsansfont{Liberation Sans}
\setmonofont{DejaVu Sans Mono}[Scale=0.84]

\usepackage{xcolor}
\definecolor{brandblue}{RGB}{20, 55, 105}
\definecolor{accentblue}{RGB}{35, 95, 165}
\definecolor{codebg}{RGB}{248, 249, 251}
\definecolor{codeframe}{RGB}{215, 222, 230}
\definecolor{javakeyword}{RGB}{0, 45, 170}
\definecolor{javastring}{RGB}{5, 120, 20}
\definecolor{javacomment}{RGB}{120, 120, 120}
\definecolor{javanumber}{RGB}{20, 75, 220}

\definecolor{noteborder}{RGB}{30, 110, 65}
\definecolor{notebg}{RGB}{242, 249, 245}
\definecolor{warnborder}{RGB}{185, 75, 10}
\definecolor{warnbg}{RGB}{254, 248, 240}
\definecolor{tipborder}{RGB}{30, 95, 160}
\definecolor{tipbg}{RGB}{242, 247, 254}

\usepackage{fancyhdr}
\usepackage{lastpage}
\pagestyle{fancy}
\fancyhf{}
\lhead{\parbox[t]{0.58\textwidth}{\raggedright\small\sffamily\color{brandblue}\textbf{T17 --- Gestione degli Errori ed Architettura delle Eccezioni}}}
\rhead{\small\sffamily\color{gray}Programmazione II -- UniVR (2025/2026)}
\cfoot{\small\sffamily Pagina \thepage\ di \pageref{LastPage}}
\renewcommand{\headrulewidth}{0.5pt}
\renewcommand{\headrule}{\hbox to\headwidth{\color{brandblue!70!black}\leaders\hrule height \headrulewidth\hfill}}

\usepackage{titlesec}
\titleformat{\section}{\Large\bfseries\sffamily\color{brandblue}}{\thesection}{0.8em}{}[\titlerule]
\titleformat{\subsection}{\large\bfseries\sffamily\color{accentblue}}{\thesubsection}{0.8em}{}
\titleformat{\subsubsection}{\normalsize\bfseries\sffamily\color{brandblue!85!black}}{\thesubsubsection}{0.8em}{}

\usepackage{needspace}
\usepackage{fancyvrb}
\DefineVerbatimEnvironment{asciibox}{Verbatim}{
  fontsize=\fontsize{7.5pt}{9.2pt}\selectfont,
  frame=single,
  rulecolor=\color{codeframe},
  fillcolor=\color{codebg},
  framesep=2.5mm
}
\usepackage{listings}
\lstset{
  backgroundcolor=\color{codebg},
  frame=single,
  rulecolor=\color{codeframe},
  basicstyle=\ttfamily\footnotesize,
  keywordstyle=\color{javakeyword}\bfseries,
  stringstyle=\color{javastring},
  commentstyle=\color{javacomment}\itshape,
  numberstyle=\tiny\color{gray},
  numbers=left,
  numbersep=7pt,
  stepnumber=1,
  breaklines=true,
  breakatwhitespace=true,
  showstringspaces=false,
  tabsize=4,
  xleftmargin=12pt,
  framexleftmargin=12pt
}

\newcommand{\passthrough}[1]{#1}
\providecommand{\tightlist}{\setlength{\itemsep}{0pt}\setlength{\parskip}{0pt}}

\usepackage{tcolorbox}
\tcbuselibrary{skins,breakable}
\newtcolorbox{studynote}[1][Nota]{
  enhanced,
  breakable,
  colback=notebg,
  colframe=noteborder,
  fonttitle=\bfseries\sffamily\small,
  coltitle=white,
  title={#1},
  arc=2mm,
  boxrule=0.8pt,
  left=9pt, right=9pt, top=7pt, bottom=7pt
}

\newtcolorbox{studywarning}[1][Attenzione]{
  enhanced,
  breakable,
  colback=warnbg,
  colframe=warnborder,
  fonttitle=\bfseries\sffamily\small,
  coltitle=white,
  title={#1},
  arc=2mm,
  boxrule=0.8pt,
  left=9pt, right=9pt, top=7pt, bottom=7pt
}

\newtcolorbox{studytip}[1][Suggerimento]{
  enhanced,
  breakable,
  colback=tipbg,
  colframe=tipborder,
  fonttitle=\bfseries\sffamily\small,
  coltitle=white,
  title={#1},
  arc=2mm,
  boxrule=0.8pt,
  left=9pt, right=9pt, top=7pt, bottom=7pt
}

\usepackage{booktabs}
\usepackage{longtable}
\usepackage{array}
\usepackage{calc}

\usepackage{hyperref}
\hypersetup{
  colorlinks=true,
  linkcolor=accentblue,
  urlcolor=accentblue,
  citecolor=brandblue,
  pdfborder={0 0 0}
}

\title{\Huge\bfseries\sffamily\color{brandblue} T17 --- Gestione degli Errori ed Architettura delle Eccezioni \\[0.3em] \large\sffamily Gerarchia Throwable, Checked vs Unchecked, try-catch-finally, Try-with-resources e AutoCloseable}
\author{\textbf{Docente:} Prof. Michele Pasqua \quad---\quad \textbf{Appunti Revisione:} Wally}
\date{A.A. 2025/2026 -- Universit\`a degli Studi di Verona}

\begin{document}
\maketitle
\thispagestyle{fancy}
\vspace{-1.2em}
\noindent\textcolor{brandblue}{\rule{\textwidth}{0.8pt}}
\vspace{1.2em}

\begin{center}\rule{0.5\linewidth}{0.5pt}\end{center}

\subsection{1. Analisi Teorica Approfondita (Slide
T17)}\label{analisi-teorica-approfondita-slide-t17}

La lezione 17 affronta la gestione degli errori e delle anomalie a
runtime attraverso il meccanismo delle \textbf{Eccezioni} in Java.

\begin{center}\rule{0.5\linewidth}{0.5pt}\end{center}

\subsubsection{1.1 Limiti della Gestione Tradizionale degli Errori
(Slide
3-6)}\label{limiti-della-gestione-tradizionale-degli-errori-slide-3-6}

Nei linguaggi procedurali privi di eccezioni (come C) o nell'uso ingenuo
dei metodi, gli errori vengono segnalati tramite \textbf{valori
sentinella} (es. restituire \passthrough{\lstinline!-1!},
\passthrough{\lstinline!null!},
\passthrough{\lstinline!Float.MAX\_VALUE!}):

\begin{lstlisting}[language=Java]
float division(int num, int den) {
    if (den != 0)
        return (float) num / den;
    else
        return Float.MAX_VALUE; // Valore speciale di errore
}
\end{lstlisting}

\paragraph{Perché questo approccio è
inadeguato?}\label{perchuxe9-questo-approccio-uxe8-inadeguato}

\begin{enumerate}
\def\labelenumi{\arabic{enumi}.}
\tightlist
\item
  \textbf{Inquinamento della logica applicativa}: Il chiamante deve
  ricordarsi di controllare continuamente i valori di ritorno con una
  ragnatela di \passthrough{\lstinline!if-else!}.
\item
  \textbf{Ambiguità semantica}: A volte il valore sentinella può essere
  un risultato matematico o di dominio assolutamente lecito.
\item
  \textbf{Inapplicabilità ai Costruttori}: Un costruttore non ha tipo di
  ritorno (\passthrough{\lstinline!void!} implicito); se i parametri
  sono non validi, non può restituire \passthrough{\lstinline!-1!} per
  segnalare il fallimento!
\item
  \textbf{Difficoltà di propagazione}: Se l'errore avviene a 10 livelli
  di profondità nello stack di chiamate, ogni singola funzione
  intermedia dovrebbe propagare manualmente il codice d'errore fino al
  punto in cui può essere gestito.
\end{enumerate}

\begin{center}\rule{0.5\linewidth}{0.5pt}\end{center}

\subsubsection{\texorpdfstring{1.2 La Filosofia delle Eccezioni in Java:
\texttt{try-catch-throw} (Slide
7-10)}{1.2 La Filosofia delle Eccezioni in Java: try-catch-throw (Slide 7-10)}}\label{la-filosofia-delle-eccezioni-in-java-try-catch-throw-slide-7-10}

Java disaccoppia la \textbf{logica di business} (\emph{cosa deve fare il
programma quando tutto va bene}) dalla \textbf{logica di gestione degli
errori} (\emph{cosa fare quando accade un'anomalia}): *
\textbf{\passthrough{\lstinline!try!}}: racchiude il blocco di codice a
rischio di eccezione. * \textbf{\passthrough{\lstinline!throw!}}:
solleva attivamente un'istanza di eccezione nel punto in cui l'anomalia
viene riscontrata. * \textbf{\passthrough{\lstinline!catch!}}: cattura
l'eccezione ed esegue il codice di ripristino/notifica. *
\textbf{Interruzione immediata}: quando viene eseguito
\passthrough{\lstinline!throw!}, il metodo \textbf{interrompe
immediatamente la propria esecuzione}, svuota lo stack frame corrente e
risale lungo la catena di chiamate (\emph{call stack}) fino al primo
blocco \passthrough{\lstinline!catch!} compatibile.

\begin{center}\rule{0.5\linewidth}{0.5pt}\end{center}

\subsubsection{1.3 La Gerarchia delle Eccezioni in Java (Slide
20-21)}\label{la-gerarchia-delle-eccezioni-in-java-slide-20-21}

Tutto l'albero discende da
\passthrough{\lstinline!java.lang.Throwable!}:

\needspace{7cm}
\begin{asciibox}
                 ┌─────────────────────┐
                 │ java.lang.Throwable │
                 └──────────▲──────────┘
                            │
        ┌───────────────────┴───────────────────┐
        │                                       │
┌───────┴─────────┐                   ┌─────────┴─────────┐
│      Error      │                   │     Exception     │
└───────▲─────────┘                   └─────────▲─────────┘
        │                                       │
┌───────┴───────────────┐               ┌───────┴───────────────┐
│  OutOfMemoryError     │               │  Checked Exceptions   │
│  StackOverflowError   │               │ (IOException, ecc.)   │
│  ThreadDeath          │               └───────────▲───────────┘
└───────────────────────┘                           │
                                        ┌───────────┴───────────┐
                                        │   RuntimeException    │
                                        │ (Unchecked Exceptions)│
                                        └───────────▲───────────┘
                                                    │
                                        ┌───────────┴───────────┐
                                        │ NullPointerException  │
                                        │ ClassCastException    │
                                        │ IndexOutOfBoundsExc.  │
                                        └───────────────────────┘
\end{asciibox}

\begin{enumerate}
\def\labelenumi{\arabic{enumi}.}
\tightlist
\item
  \textbf{\passthrough{\lstinline!Error!} (Unchecked)}:

  \begin{itemize}
  \tightlist
  \item
    Condizioni anomale e catastrofiche interne alla JVM o al sistema
    operativo (\passthrough{\lstinline!OutOfMemoryError!},
    \passthrough{\lstinline!StackOverflowError!}).
  \item
    Il codice utente non deve tentare di catturarle, poiché lo stato
    della memoria della JVM è compromesso.
  \end{itemize}
\item
  \textbf{\passthrough{\lstinline!Exception!}}:

  \begin{itemize}
  \tightlist
  \item
    Anomalie derivanti da condizioni esterne o logiche. Si ramificano in
    due grandi famiglie:

    \begin{itemize}
    \tightlist
    \item
      \textbf{Checked Exceptions} (sottoclassi dirette di
      \passthrough{\lstinline!Exception!}, escluse le
      \passthrough{\lstinline!RuntimeException!}).
    \item
      \textbf{Unchecked Exceptions} (sottoclassi di
      \passthrough{\lstinline!RuntimeException!}).
    \end{itemize}
  \end{itemize}
\end{enumerate}

\begin{center}\rule{0.5\linewidth}{0.5pt}\end{center}

\subsubsection{1.4 Checked vs Unchecked Exceptions (Slide
22-24)}\label{checked-vs-unchecked-exceptions-slide-22-24}

\begin{longtable}[]{@{}
  >{\raggedright\arraybackslash}p{(\linewidth - 4\tabcolsep) * \real{0.3333}}
  >{\raggedright\arraybackslash}p{(\linewidth - 4\tabcolsep) * \real{0.3333}}
  >{\raggedright\arraybackslash}p{(\linewidth - 4\tabcolsep) * \real{0.3333}}@{}}
\toprule\noalign{}
\begin{minipage}[b]{\linewidth}\raggedright
Caratteristica
\end{minipage} & \begin{minipage}[b]{\linewidth}\raggedright
Checked Exceptions (\passthrough{\lstinline!extends Exception!})
\end{minipage} & \begin{minipage}[b]{\linewidth}\raggedright
Unchecked Exceptions
(\passthrough{\lstinline!extends RuntimeException!})
\end{minipage} \\
\midrule\noalign{}
\endhead
\bottomrule\noalign{}
\endlastfoot
\textbf{Natura dell'errore} & Situazioni anomale \textbf{prevedibili} ma
esterne (es. \passthrough{\lstinline!FileNotFoundException!}, rete
disconnessa). & \textbf{Errori del programmatore} (bug logici,
precondizioni violate, puntatori nulli). \\
\textbf{Controllo del compilatore} & \textbf{Obbligatorio} (\emph{Catch
or Declare rule}). Il codice non compila se non gestite con
\passthrough{\lstinline!try-catch!} o dichiarate con
\passthrough{\lstinline!throws!}. & \textbf{Opzionale}. Il compilatore
non richiede alcuna dichiarazione o cattura. \\
\textbf{Firma del metodo} & Modificata:
\passthrough{\lstinline!public void foo() throws MyCheckedException!} &
Invariata: nessun \passthrough{\lstinline!throws!} richiesto. \\
\textbf{Diffusione (\emph{Propagazione})} & \emph{``Si propagano come un
virus''} (Slide 24): costringono tutti i chiamanti intermedi a
dichiarare \passthrough{\lstinline!throws!} o gestire. & Risalgono
naturalmente lo stack fino al \passthrough{\lstinline!main!} o al
gestore globale. \\
\textbf{Raccomandazione moderna (Slide 24)} & Definirle solo se l'utente
del metodo \textbf{può ragionevolmente intraprendere un'azione di
recupero}. & \textbf{Preferite} per errori di validazione dello stato o
degli argomenti. \\
\end{longtable}

\begin{center}\rule{0.5\linewidth}{0.5pt}\end{center}

\subsubsection{1.5 ``Exception Dirty Tricks'': Trasformare Checked in
Unchecked (Slide
25-26)}\label{exception-dirty-tricks-trasformare-checked-in-unchecked-slide-25-26}

Un pattern comune quando una checked exception sporcherebbe decine di
firme senza che il chiamante intermedio possa gestirla:

\begin{lstlisting}[language=Java]
public void bar() { // Nessuna dichiarazione throws richiesta!
    try {
        foo(); // foo() dichiara throws MyCheckedException
    } catch (MyCheckedException e) {
        throw new RuntimeException(e); // Incapsulata e rilanciata come Unchecked!
    }
}
\end{lstlisting}

L'eccezione originale non viene persa: è memorizzata come causa
(\emph{chained exception} visibile nel messaggio
\passthrough{\lstinline!Caused by:!} dello stacktrace).

\begin{center}\rule{0.5\linewidth}{0.5pt}\end{center}

\subsubsection{\texorpdfstring{1.6 Eccezioni nei Cicli (\emph{Exceptions
and Loops}, Slide
27-28)}{1.6 Eccezioni nei Cicli (Exceptions and Loops, Slide 27-28)}}\label{eccezioni-nei-cicli-exceptions-and-loops-slide-27-28}

\begin{itemize}
\tightlist
\item
  \textbf{Errore sulla singola iterazione (Slide 27)}: il blocco
  \passthrough{\lstinline!try-catch!} è \textbf{all'interno} del ciclo.
  Se l'iterazione fallisce, viene intercettata e il ciclo prosegue con
  la successiva (es. lettura da input utente fino a valore corretto).
\item
  \textbf{Errore fatale per l'intero ciclo (Slide 28)}: il blocco
  \passthrough{\lstinline!try-catch!} \textbf{avvolge} il ciclo. Se si
  verifica un'anomalia, il ciclo viene interrotto definitivamente e
  l'esecuzione salta direttamente al \passthrough{\lstinline!catch!}
  esterno.
\end{itemize}

\begin{center}\rule{0.5\linewidth}{0.5pt}\end{center}

\subsection{\texorpdfstring{2. Analisi Dettagliata del Codice
(\texttt{Lezione16/MavenDate})}{2. Analisi Dettagliata del Codice (Lezione16/MavenDate)}}\label{analisi-dettagliata-del-codice-lezione16mavendate}

In \passthrough{\lstinline!Lezione16/MavenDate!} il prof. Pasqua applica
i principi della Slide T17 rimuovendo i vecchi
\passthrough{\lstinline'System.out.println("Illegal date!")'} e
introducendo la gestione strutturata degli errori tramite il package
\passthrough{\lstinline!it.oop.exception!}.

\subsubsection{\texorpdfstring{2.1 Le Nuove Eccezioni:
\href{file:///home/wally/Documenti/GitHub/Programmazione-II/25-26/Teoria-e-Esercitazioni/Codice-20260902/Lezione16/MavenDate/src/main/java/it/oop/exception/IllegalDateException.java}{IllegalDateException.java}
e
\href{file:///home/wally/Documenti/GitHub/Programmazione-II/25-26/Teoria-e-Esercitazioni/Codice-20260902/Lezione16/MavenDate/src/main/java/it/oop/exception/OrderdPairException.java}{OrderdPairException.java}}{2.1 Le Nuove Eccezioni: IllegalDateException.java e OrderdPairException.java}}\label{le-nuove-eccezioni-illegaldateexception.java-e-orderdpairexception.java}

\paragraph{\texorpdfstring{1.
\href{file:///home/wally/Documenti/GitHub/Programmazione-II/25-26/Teoria-e-Esercitazioni/Codice-20260902/Lezione16/MavenDate/src/main/java/it/oop/exception/IllegalDateException.java}{IllegalDateException.java}
(Unchecked)}{1. IllegalDateException.java (Unchecked)}}\label{illegaldateexception.java-unchecked}

\begin{lstlisting}[language=Java]
package it.oop.exception;

public class IllegalDateException extends RuntimeException {
    public IllegalDateException(String message) {
        super(message);
    }
}
\end{lstlisting}

\begin{itemize}
\tightlist
\item
  Estende \passthrough{\lstinline!RuntimeException!}: è una
  \textbf{Unchecked Exception}.
\item
  Modella una violazione delle precondizioni sui valori numerici
  (giorno, mese, anno).
\item
  Non costringe tutti i client a dichiarare
  \passthrough{\lstinline!throws IllegalDateException!}.
\end{itemize}

\paragraph{\texorpdfstring{2.
\href{file:///home/wally/Documenti/GitHub/Programmazione-II/25-26/Teoria-e-Esercitazioni/Codice-20260902/Lezione16/MavenDate/src/main/java/it/oop/exception/OrderdPairException.java}{OrderdPairException.java}
(Checked)}{2. OrderdPairException.java (Checked)}}\label{orderdpairexception.java-checked}

\begin{lstlisting}[language=Java]
package it.oop.exception;

public class OrderdPairException extends Exception { // Notare il typo del docente "Orderd"
    public OrderdPairException(String message) {
        super(message);
    }
}
\end{lstlisting}

\begin{itemize}
\tightlist
\item
  Estende direttamente \passthrough{\lstinline!Exception!}: è una
  \textbf{Checked Exception}.
\item
  \textbf{Impatto immediato sulla firma}:

  \begin{itemize}
  \item
    In
    \href{file:///home/wally/Documenti/GitHub/Programmazione-II/25-26/Teoria-e-Esercitazioni/Codice-20260902/Lezione16/MavenDate/src/main/java/it/oop/core/OrderedPair.java\#L9}{OrderedPair.java}:

\begin{lstlisting}[language=Java]
public OrderedPair(T first, T second) throws OrderdPairException {
    this.first = first;
    this.second = second;
    if (first.compareTo(second) > 0)
        throw new OrderdPairException("Not orderd pair");
}
\end{lstlisting}
  \item
    In
    \href{file:///home/wally/Documenti/GitHub/Programmazione-II/25-26/Teoria-e-Esercitazioni/Codice-20260902/Lezione16/MavenDate/src/main/java/it/oop/core/DateInterval.java\#L6}{DateInterval.java}:

\begin{lstlisting}[language=Java]
public DateInterval(Date left, Date right) throws OrderdPairException {
    super(left, right); // Invoca super() che lancia la checked exception, quindi DateInterval DEVE dichiarare throws!
}
\end{lstlisting}
  \end{itemize}
\end{itemize}

\begin{center}\rule{0.5\linewidth}{0.5pt}\end{center}

\subsubsection{\texorpdfstring{2.2 Refactoring del Metodo
\texttt{verify()} in
\href{file:///home/wally/Documenti/GitHub/Programmazione-II/25-26/Teoria-e-Esercitazioni/Codice-20260902/Lezione16/MavenDate/src/main/java/it/oop/core/Date.java\#L28-L35}{Date.java}}{2.2 Refactoring del Metodo verify() in Date.java}}\label{refactoring-del-metodo-verify-in-date.java}

\begin{lstlisting}[language=Java]
    void verify() {
        if (year < 0)
            throw new IllegalDateException("Illegal date: wrong year");
        if (month < 1 || month > 12)
            throw new IllegalDateException("Illegal date: wrong month");
        if (day < 1 || day > daysPerMonth(month))
            throw new IllegalDateException("Illegal date: wrong day");
    }
\end{lstlisting}

\begin{itemize}
\tightlist
\item
  \textbf{Miglioramento architetturale}:

  \begin{itemize}
  \tightlist
  \item
    Prima: un giorno errato (es. 35) stampava un messaggio sulla console
    ma creava comunque un oggetto \passthrough{\lstinline!Date!}
    incoerente nell'Heap.
  \item
    Ora: il costruttore viene interrotto da
    \passthrough{\lstinline!throw!} prima che l'assegnazione sia
    completata. \textbf{È impossibile creare un'istanza con stato non
    valido!}
  \end{itemize}
\end{itemize}

\begin{center}\rule{0.5\linewidth}{0.5pt}\end{center}

\subsubsection{\texorpdfstring{2.3
\href{file:///home/wally/Documenti/GitHub/Programmazione-II/25-26/Teoria-e-Esercitazioni/Codice-20260902/Lezione16/MavenDate/src/main/java/it/oop/ui/MainDate.java}{MainDate.java}
e l'Applicazione di Tutti i Pattern della Slide
T17}{2.3 MainDate.java e l'Applicazione di Tutti i Pattern della Slide T17}}\label{maindate.java-e-lapplicazione-di-tutti-i-pattern-della-slide-t17}

\begin{lstlisting}[language=Java]
package it.oop.ui;

import it.oop.core.*;
import it.oop.core.Date;
import it.oop.exception.IllegalDateException;
import it.oop.exception.OrderdPairException;

import java.util.InputMismatchException;
import java.util.Scanner;

public class MainDate {

    public static void main(String[] args) {
        System.out.println("Insert day, month, year as numbers");
        int d = -1, m = -1, y = -1;
        boolean correct = false;

        // 1. Pattern Slide 27: try-catch dentro il ciclo per recuperare l'errore di input dell'utente
        while (correct == false) {
            try {
                Scanner sc = new Scanner(System.in);
                d = sc.nextInt();
                m = sc.nextInt();
                y = sc.nextInt();
                correct = true; // Se uno dei nextInt() fallisce, questa riga non viene raggiunta
            } catch (InputMismatchException ime) {
                System.out.println("Input not valid, retry");
            }
        }

        // 2. Cattura della Unchecked Exception IllegalDateException per notifica pulita
        try {
            FormattedDate date = new ItalianDate(d, m, y);
            System.out.println(date.toString());
        } catch (IllegalDateException ide) {
            System.out.println("Date not valid: " + ide.getMessage());
        }

        // 3. Pattern Slide 26 ("Dirty Trick"): incapsulamento di Checked Exception in RuntimeException
        try {
            DateInterval di = new DateInterval(new Date(31, 1, 2025), new Date(1, 1, 2025));
        } catch (OrderdPairException ope) {
            ope.printStackTrace();
            throw new RuntimeException(ope); // Rilanciata come unchecked!
        }
    }
}
\end{lstlisting}

\begin{center}\rule{0.5\linewidth}{0.5pt}\end{center}

\subsection{3. Risultato di Esecuzione Reale da
Terminale}\label{risultato-di-esecuzione-reale-da-terminale}

\subsubsection{\texorpdfstring{Test 1: Input con data errata
(\texttt{35\ 1\ 2025})}{Test 1: Input con data errata (35 1 2025)}}\label{test-1-input-con-data-errata-35-1-2025}

Esecuzione:

\begin{lstlisting}[language=bash]
mvn exec:java -Dexec.mainClass="it.oop.ui.MainDate" <<< "35 1 2025"
\end{lstlisting}

Output:

\begin{lstlisting}
Insert day, month, year as numbers
Date not valid: Illegal date: wrong day
it.oop.exception.OrderdPairException: Not orderd pair
    at it.oop.core.OrderedPair.<init>(OrderedPair.java:13)
    at it.oop.core.DateInterval.<init>(DateInterval.java:7)
    at it.oop.ui.MainDate.main(MainDate.java:35)
...
Caused by: it.oop.exception.OrderdPairException: Not orderd pair
\end{lstlisting}

\subsubsection{Analisi dell'output:}\label{analisi-delloutput}

\begin{enumerate}
\def\labelenumi{\arabic{enumi}.}
\tightlist
\item
  \passthrough{\lstinline!Date not valid: Illegal date: wrong day!}: il
  costruttore di \passthrough{\lstinline!ItalianDate!} delega a
  \passthrough{\lstinline!Date.verify()!}, che lancia
  \passthrough{\lstinline!IllegalDateException("Illegal date: wrong day")!},
  catturata dal blocco
  \passthrough{\lstinline!catch (IllegalDateException ide)!} che stampa
  il messaggio di errore controllato.
\item
  \passthrough{\lstinline!DateInterval di = new DateInterval(31/1/2025, 1/1/2025)!}:
  poiché la data iniziale \passthrough{\lstinline!31/1/2025!} è
  cronologicamente successiva alla data finale
  \passthrough{\lstinline!1/1/2025!} (violazione dell'invariante
  \passthrough{\lstinline!start <= end!}),
  \passthrough{\lstinline!OrderedPair!} solleva la checked
  \passthrough{\lstinline!OrderdPairException!}. Il blocco catch stampa
  lo stack trace con \passthrough{\lstinline!ope.printStackTrace()!} e
  poi rilancia l'errore incapsulato con
  \passthrough{\lstinline!throw new RuntimeException(ope);!}, terminando
  l'esecuzione con codice di uscita d'errore (build failure) come
  previsto.
\end{enumerate}

\begin{center}\rule{0.5\linewidth}{0.5pt}\end{center}

Dimmi \textbf{``vai''} per procedere con il \textbf{Blocco 17} (Lezione
18 / Slide T18 --- \emph{Streams and Functional Programming}, Stream
API, operazioni intermedie vs terminali, \passthrough{\lstinline!map!},
\passthrough{\lstinline!filter!}, \passthrough{\lstinline!reduce!},
\passthrough{\lstinline!collect!} e il nuovo file
\passthrough{\lstinline!MainStream.java!} in
\passthrough{\lstinline!Lezione17/MavenDate!}).


\end{document}
